\section{Related Work}\label{sec:related}

\paragraph{Work on embedding compression}
Work on embedding compression decides which IDs share parameters. The hashing trick maps IDs to rows uniformly~\cite{weinberger2009hashing}, and hash embeddings combine several hashed rows with learned importance weights~\cite{svenstrup2017hashemb}. Frequency-based double hashing gives popular items private rows and hashes the rest~\cite{zhang2020doublehash}. Later methods learn or adapt the assignment: Learning to Collide trains the hash function~\cite{ghaemmaghami2022collide}, CCE re-clusters shared rows during training~\cite{tsang2023cce}, and CAFE allocates exclusive rows to IDs whose importance it estimates online~\cite{zhang2024cafe}. Graph-based methods assign by interaction structure: GraphHash clusters the user--item graph by modularity~\cite{wu2025graphhash,blondel2008louvain}, BACO balances co-clusters of users and items~\cite{jiang2026baco}, and GDE replaces hard buckets with soft gating~\cite{lee2026gde}, so items under GDE do not tie exactly. These works do not ask what sharing does to embedding magnitude or to the popularity signal it carries; they compare assignments under a single training objective.

\paragraph{Work on embedding magnitude}
Work on embedding magnitude shows that, on full tables, the norm of an item embedding encodes its popularity. PRISM traces this to weight decay, which shrinks rarely updated items~\cite{loveland2025prism}; Adap-$\tau$ adapts the temperature to control magnitude~\cite{chen2023adaptau}; DDC decomposes embeddings analytically and isolates a popularity direction~\cite{liu2026ddc}; and DNA documents the same norm bias in GCN-based recommenders~\cite{choi2026dna}. Training objectives differ in whether they see this signal: BPR scores by dot product~\cite{rendle2009bpr}, whereas DirectAU optimises alignment and uniformity of $\ell_2$-normalised embeddings~\cite{wang2022directau}. We build on the finding that magnitude encodes popularity. These works do not study shared tables, where one row's norm must serve every item in its bucket.

\paragraph{Work on popularity terms}
Work on popularity terms adds item-level information to the score, either learned during training or fixed at score time. MEmCom attaches a learned per-item scale to a shared row~\cite{pansare2022memcom}. Logit adjustment~\cite{menon2021logit} and the logQ correction~\cite{yi2019logq} add log-frequency terms to the logits, and causal approaches remove or re-inject popularity at inference~\cite{zhang2021pda,wei2021macr}. Our \crs{} is the log-prior score form of Cheng~\cite{cheng2026rankone} (eq.~53), applied after training with one validation-tuned weight. These works do not connect such terms to what sharing erases, nor show why a learned per-item scale cannot carry popularity under DirectAU (\S\ref{sec:theory}).
